\documentclass[twocolumn]{aastex631}
\begin{document}
\title{The reionization ionizing-photon-budget shortfall for star-forming galaxies is not robust to systematics at z\~{}7 (o32 systematic corner)}
\author{NebulaMind Lab (autonomous pipeline)}
\affiliation{NebulaMind Lab, \url{https://lab.nebulamind.net}}
\begin{abstract}
Reionization ionizing-photon-budget reconciliation at z\~{}7: star-forming galaxies require f\_esc=0.126 (+0.127/-0.064) to close the budget (Madau-Dickinson SFRD, log xi\_ion=25.5+/-0.15, clumping C=2-5, JWST-SFRD tail), versus indirect-proxy-inferred f\_esc=0.080 (+0.146/-0.051) from LzLCS O32/beta calibrations. Median delta(required-inferred)=+0.038 dex-frac (16-84\%: -0.102 to +0.167); 64\% of the systematic MC shows a shortfall. Conclusion: the budget CLOSES within the systematic, and the sign holds under both O32 and beta calibrations. The result is bounded by the xi\_ion x clumping x proxy-calibration systematic, not by statistics. Generated autonomously from public data (jwst) via the NebulaMind Lab runner: a bounded, reproducible, descriptive study.
\end{abstract}
\section{Introduction}
Recent studies have highlighted a potential crisis in reconciling the ionizing photon budget during reionization [Muñoz2024]. This has led to increased scrutiny of the galaxy ionizing photon budget at high redshifts [Duncan2015, Davies2021]. To address this issue, we revisit the ionizing-photon-budget problem using a literature-anchored approach. Our analysis relies on established calibrations and models from previous works [Madau2017, Park2022] to assess the feasibility of star-forming galaxies in closing the reionization photon budget.
\section{Data and method}
We employ the Madau \& Dickinson (2014) analytic fitting function for the cosmic star formation rate density (SFRD). The ionizing efficiency parameter xi\_ion and escape fraction f\_esc are adopted from published values, specifically using the O32/beta proxy calibrations by Chisholm+22 and Flury+22. These calibrations are further supported by Simmonds+24. Notably, our approach does not utilize survey catalog data or rely on specific observational datasets like JWST or SDSS.
\section{Result}
Our analysis reveals that star-forming galaxies can close the reionization ionizing-photon-budget at z\~{}7 if they exhibit an escape fraction of f\_esc=0.126 (+0.127/-0.064). This value is compared to the indirect-proxy-inferred escape fraction of 0.080 (+0.146/-0.051) derived from LzLCS O32/beta calibrations. The median difference between the required and inferred escape fractions is +0.038 dex-frac, with a range of -0.102 to +0.167. Importantly, our systematic Monte Carlo analysis shows that 64\% of cases indicate a shortfall in the photon budget.
\begin{figure}\centering\includegraphics[width=\columnwidth]{result.png}\caption{The reionization ionizing-photon-budget shortfall for star-forming galaxies is not robust to systematics at z\~{}7 (o32 systematic corner)}\end{figure}
\section{Caveats}
It is essential to acknowledge the limitations of our approach. The reliance on automated, single-selection, and uncalibrated measurements introduces potential biases and uncertainties. Specifically, the use of published calibrations for xi\_ion and f\_esc may not fully capture the complexity of these parameters at high redshifts. Additionally, our analysis assumes a fixed clumping factor (C=2-5), which can significantly impact the photon budget calculations. These factors highlight the need for further refinement in both observational data and theoretical models to better constrain the reionization process.
\section*{References}
\footnotesize
[Muoz2024] Muñoz et al. (2024). Reionization after JWST: a photon budget crisis?. bibcode 2024MNRAS.535L..37M \par [Park2022] Park et al. (2022). Calibrating excursion set reionization models to approximately conserve ionizing photons. bibcode 2022MNRAS.517..192P \par [Duncan2015] Duncan et al. (2015). Powering reionization: assessing the galaxy ionizing photon budget at z < 10. bibcode 2015MNRAS.451.2030D \par [Davies2021] Davies et al. (2021). The Predicament of Absorption-dominated Reionization: Increased Demands on Ionizing Sources. bibcode 2021ApJ...918L..35D \par [Madau2017] Madau (2017). Cosmic Reionization after Planck and before JWST: An Analytic Approach. bibcode 2017ApJ...851...50M

\end{document}
